\section{Reproduction}
\label{sec:reproduction}

All commands are run from the \garnet{} repository root on a merlin7 login
node, unless stated otherwise.

\paragraph{1. Build \ffn{}.} The tarballs \code{petsc-3.19.6}, \code{slepc-3.19.2}
and \code{dealii-9.5.2} must be in
\code{/data/scratch/shared/poncet\_m/femffusion-deps-v1/src}. Then:
\begin{verbatim}
sbatch slurm/build_femffusion.sh
# or, on a login node: BUILD_JOBS=6 bash -l slurm/build_femffusion.sh
\end{verbatim}
This builds the libraries and \code{\textasciitilde/FEMFFUSION/build-release/femffusion.exe},
writes \code{build-info.json} (versions, flags and the runtime library path)
and runs \code{femffusion.exe -t}.

\paragraph{2. HP-MR input for \ffn{}.}
\begin{verbatim}
PYTHONPATH=. python tools/export_femffusion.py \
    /data/scratch/shared/poncet_m/garnet-femffusion-perf/inputs/hpmr-v2
\end{verbatim}
The output is the hexagonal mesh, the XML materials (kinetics included), the
steady \code{.prm} and the transient \code{.prm} with its final-state XML. The
steady runs used \code{hpmr-v1}, written before the kinetics moved into the
XML. Its static data are identical, and both inputs give the same $k$.

\paragraph{3. Runs.} Every run is one line of a plan file in
\code{benchmark-results/garnet-vs-femffusion/plans/}:
\code{<chain> <tag> <tool> <arguments>}, where the tool is
\code{tools/assess\_garnet\_vs\_femffusion.py} or \code{tools/run\_femffusion.py}.
\begin{verbatim}
bash slurm/stage_and_submit_femffusion_cmp.sh cpu-hourly plans/cpu-steady-a.plan
bash slurm/stage_and_submit_femffusion_cmp.sh cpu-daily  plans/cpu-steady-b.plan
bash slurm/stage_and_submit_femffusion_cmp.sh gpu        plans/gpu-steady.plan
\end{verbatim}
Each submission stages its own snapshot of the repository with a
\code{PROVENANCE.txt} (\garnet{} commit, status and diff; \ffn{} commit and
build). \code{slurm/run\_perf\_plan.sh} runs the chains, pinning each chain on
a CPU node to its own NUMA domain. The plans used, in order:
\begin{itemize}
\item tuning: \code{cpu-ff-tune}, \code{cpu-ff-trtune}, \code{cpu-garnet-tune},
  \code{gpu-garnet-tune};
\item references: \code{gpu-ref}, \code{cpu-ff-ref}, and \code{cpu-ref-shared}
  (target \code{cpu-daily-shared}: part of a shared node, for runs whose time
  is not used);
\item steady: \code{gpu-steady}, \code{cpu-steady-a/b}, \code{cpu-alone};
\item transient: \code{gpu-transient}, \code{gpu-transient-2},
  \code{cpu-transient-a/b}, \code{cpu-fftransient-a/b/c};
\item runs left over by the queue limits: \code{cpu-split-h1/h2/h3},
  \code{cpu-split-d} and the single-run plans \code{cpu-one-*}, each an
  exclusive node with a short time limit.
\item Wielandt shift, HP-MR steady: \code{gpu-wielandt} (\code{gpu-hourly})
  and \code{cpu-wielandt} (\code{cpu-hourly}). Both run the fixed and the
  adaptive shift; \code{WIELANDT\_MODE} in \code{make\_figures.py} selects
  the rows. The \code{tune-floor-*} runs of \code{cpu-wielandt} are the
  iteration-error check of Section~\ref{sec:method}. \code{cpu-one-hpmr-wf-r8}
  replaced the power-iteration plan \code{cpu-one-hpmr-r8}, cancelled
  before it ran. \code{gpu-2g-wielandt} (\code{gpu-hourly}) checks the two
  two-group cores at Langenbuch $r=12$ and IAEA $r=24$ (tag \code{tune-},
  not rows).
\item IAEA-3D fine meshes: \code{gpu-iaea-fine} (\code{gpu-hourly}),
  $r$ = 20, 24, 28, with an $r=16$ control (tag \code{tune-}, not a row)
  that reproduces the \code{gpu-steady} $r=16$ run bit for bit on the
  current code.
\end{itemize}

\paragraph{4. Report.}
\begin{verbatim}
cd benchmark-results/garnet-vs-femffusion
bash collect.sh                      # results -> data/perf/<job>/
PYTHONPATH=../.. python make_reference.py
make PYTHON=python                   # figures, tables, data/summary.json, PDF
\end{verbatim}

\paragraph{Checks in the test suite.}
\begin{itemize}
\item \code{tests/validation/test\_femffusion\_equivalence.py}: \garnet's
  IAEA-3D and Langenbuch builders pose \ffn's example inputs, node by node.
\item \code{tests/verification/test\_femffusion\_io.py}: the XML writer round
  trip and its kinetics layout.
\item \code{tests/validation/test\_hpmr.py}: the hex-homogenized drum option
  (\code{absorber="smeared"}).
\end{itemize}
